\documentclass[a4paper]{article}
% Español (México) con XeLaTeX: fontspec para fuentes Unicode, polyglossia
% para la división silábica y los nombres del documento en la variante mexicana.
\usepackage{fontspec}
\usepackage{polyglossia}
\setdefaultlanguage[variant=mexican]{spanish}
% Los nombres chinos van como «pinyin (original)»: xeCJK da a los caracteres
% Han su propia fuente; sin ella XeLaTeX los omite con un aviso.
% FandolSong no cubre algunos caracteres de nombres (U+73FA); WenQuanYi Zen Hei sí.
\usepackage[AutoFallBack=true]{xeCJK}
\setCJKmainfont{FandolSong-Regular.otf}
\setCJKfallbackfamilyfont{\CJKrmdefault}{WenQuanYi Zen Hei}
% polyglossia no traduce los nombres de listings.
\AtBeginDocument{\providecommand{\lstlistingname}{}\renewcommand{\lstlistingname}{Listado}%
  \providecommand{\lstlistlistingname}{}\renewcommand{\lstlistlistingname}{Índice de listados}}
\usepackage{amsmath, amssymb}
\usepackage{graphicx}
\usepackage[margin=2.35cm]{geometry}
\usepackage[most]{tcolorbox}
\usepackage{booktabs}
\usepackage{tabularx}
\usepackage{longtable}
\usepackage{array}
\usepackage{float}
\usepackage{xcolor}
\usepackage{hyperref}
\usepackage{fancyhdr}
\hypersetup{colorlinks=true, linkcolor=blue!55!black, urlcolor=blue!60!black}
\graphicspath{{./}{figures/}}
\setlength{\parskip}{0.55em}
\setlength{\parindent}{2em}
\setlength{\emergencystretch}{3em}
\renewcommand{\arraystretch}{1.18}
\newtcolorbox{knowledgebox}[1]{enhanced, colback=blue!5!white, colframe=blue!70!black, colbacktitle=blue!70!black, coltitle=white, fonttitle=\bfseries, title={#1}, sharp corners, breakable}
\newtcolorbox{importantbox}[1]{enhanced, colback=yellow!10!white, colframe=yellow!75!black, colbacktitle=yellow!75!black, coltitle=black, fonttitle=\bfseries, title={#1}, sharp corners, breakable}
\newtcolorbox{warningbox}[1]{enhanced, colback=red!5!white, colframe=red!70!black, colbacktitle=red!70!black, coltitle=white, fonttitle=\bfseries, title={#1}, sharp corners, breakable}
\pagestyle{fancy}
\fancyhf{}
\fancyfoot[C]{\thepage}
\fancyfoot[R]{\footnotesize\color{gray}\href{https://github.com/hqhq1025/ai-course-notes}{AI Course Notes}}
\renewcommand{\headrulewidth}{0pt}
\renewcommand{\footrulewidth}{0pt}
\newcommand{\videourl}{https://www.youtube.com/watch?v=iiBY0fqpThI}
\newcommand{\repourl}{https://github.com/hqhq1025/ai-course-notes}

\begin{document}
\begin{titlepage}
\centering
{\Large Notas didácticas de una entrevista a fondo\par}
\vspace{1.0cm}
{\huge\bfseries Modelos del mundo, criterio de investigación y «escapar de Silicon Valley»\par}
\vspace{0.5cm}
{\Large Zhang Xiaojun conversa con Xie Saining (谢赛宁)\par}
\vspace{0.35cm}
{\large Elaboradas a partir del video público de Zhang Xiaojun Podcast, la transcripción en GPU, la estructura de capítulos y diagramas conceptuales propios\par}
\vspace{0.3cm}
{\large 2026-03-16\par}
\vspace{0.85cm}
\includegraphics[width=0.88\textwidth,height=0.42\textheight,keepaspectratio]{cover.jpg}\par
\vfill
\begin{tcolorbox}[width=0.92\textwidth, colback=black!2!white, colframe=black!55, sharp corners]
\textbf{Título del video}: 133. Entrevista maratónica de 7 horas con Xie Saining: modelos del mundo, escapar de Silicon Valley, AMI Labs, dos rechazos a Ilya, Yann LeCun (杨立昆), Fei-Fei Li (李飞飞) y 42\par
\textbf{Canal del video}: Zhang Xiaojun Podcast\par
\textbf{Duración del video}: 06:45:29\par
\textbf{Enlace del video}: \href{\videourl}{\nolinkurl{\videourl}}\par
\textbf{Nota sobre el material}: YouTube no ofrece subtítulos disponibles; el texto se elaboró a partir de la transcripción con faster-whisper large-v3 en CUDA, los capítulos del video, la portada y diagramas didácticos propios. Las tomas fijas de la entrevista no se repiten en el cuerpo del texto.\par
\textbf{Página del proyecto}: \href{\repourl}{\nolinkurl{\repourl}}\par
\end{tcolorbox}
\end{titlepage}

\tableofcontents
\newpage
\section{Introducción: no es un perfil biográfico, sino una línea de investigación}

Esta entrevista dura seis horas y cuarenta y cinco minutos. En apariencia, es la primera vez que Xie Saining (谢赛宁) relata de forma sistemática su formación, sus estudios, su investigación, su emprendimiento y sus criterios de valor. El hilo más profundo es otro: cómo un investigador que trabaja en computer vision y representation learning redefine «visión», «modelo del mundo» y «cerebro robótico» en un momento en que la narrativa de los LLM se impone sobre todo lo demás. Por eso, esta nota no repetirá la entrevista palabra por palabra como un relato cronológico, sino que la organiza como una línea de investigación que se puede estudiar.

Este episodio tiene tres ejes. El primero es la trayectoria personal: de la clase ACM de la Universidad Jiao Tong de Shanghái a UCSD, FAIR, DeepMind, NYU y AMI Labs; cada decisión giró en torno a «con quién trabajar, qué problema abordar y si hay recursos suficientes». El segundo es la línea técnica: de las tareas de visión, las representaciones jerárquicas y la comprensión de video al predictive world model. El tercero es la línea organizacional: cuando un problema de investigación requiere recursos muy superiores a los de un laboratorio orientado a publicar artículos, pero no conviene que la carrera armamentista de productos lo absorba por completo, ¿cómo puede una empresa emergente conservar el research oxygen?

\begin{importantbox}{Tesis central del episodio}
Xie Saining insiste a lo largo de la entrevista en que los LLM son un componente importante de los sistemas inteligentes, pero no son todo. Una inteligencia orientada de verdad al mundo físico necesita aprender representaciones relevantes para la tarea a partir de señales continuas, de alta dimensión y con ruido, predecir las consecuencias de una action y usar esas predicciones para planning, decision making y control de seguridad.
\end{importantbox}

\begin{knowledgebox}{Nota sobre la estrategia visual}
Este video muestra una toma fija de entrevista, sin slides didácticas, pizarrón, demostraciones de producto ni gráficas legibles. Conforme al estándar de pódcast de este repositorio, el cuerpo del texto no repite fotogramas de los participantes; la portada sirve para identificar la fuente, y el contenido didáctico se presenta mediante el ciclo del modelo del mundo, la escalera de capacidades, el modelo organizacional y tablas de conceptos.
\end{knowledgebox}

\subsection{Resumen de la sección}

El valor de este episodio no está en saber a quién rechazó Xie Saining o dónde trabajó, sino en entender cómo una postura técnica se forma poco a poco a partir de la experiencia personal, el criterio de investigación y las condiciones organizacionales. Las secciones siguientes dividen la entrevista en: la trayectoria de investigación, la definición de modelo del mundo, el límite entre los LLM y el mundo físico, la propuesta organizacional de AMI y, al final, la reflexión sobre la inteligencia y las citas filosóficas.

\section{The Normal One: las decisiones de investigación detrás del relato de una persona común}

La entrevista comienza hablando de Nueva York, del podcast, de la infancia y de la Universidad Jiao Tong de Shanghái (上海交大). Xie Saining (谢赛宁) repite que él no es the chosen one, sino the normal one. Es fácil malinterpretar esta frase como un gesto de modestia; leída en el contexto de toda la entrevista, se parece más a la forma en que un investigador se sitúa a sí mismo: no hay que explicar el éxito como un guion de genio otorgado por el destino, sino como el resultado de la acción conjunta de decisiones sostenidas en el tiempo, personas y organizaciones, curiosidad, fracasos y azar.

Cuando habla de su infancia, aparecen dos elementos una y otra vez: primero, su madre lo llevaba a viajar por muchos lugares; segundo, los numerosos libros del estudio de su padre. Los viajes lo pusieron en contacto con la diversidad del mundo real, y la lectura le dio una entrada al mundo abstracto. Cuando a los nueve años tuvo computadora y conoció los videojuegos e internet, sintió por primera vez la explosión del «contenido» y de la expresión. En un perfil biográfico estos detalles son una historia de formación; en una nota técnica corresponden a las dos caras de los modelos del mundo que se tratan más adelante: al mundo hay que entrar en persona, y la abstracción requiere un modelado continuo.

\begin{knowledgebox}{Leer la línea técnica a partir de la trayectoria personal}
En esta entrevista, la experiencia personal no es una anécdota de fondo. Los viajes, los libros, los videojuegos, internet, el cine, las calles de Nueva York, las estancias de investigación y la organización de una empresa emergente responden todos a la misma pregunta: ¿de dónde debería aprender el mundo un sistema inteligente? Si aprende solo del texto, obtiene el relato que los humanos ya comprimieron; si ha de entender el mundo físico, necesita enfrentarse de nuevo a señales continuas, acciones corporales, retroalimentación del entorno y la vida real.
\end{knowledgebox}

\subsection{La clase ACM, la formación general y el rechazo a la competencia excesiva}

Lo anterior sobre la infancia y la entrada a internet sirve para mostrar que el interés nunca surge de un solo punto; esta sección pasa a la clase ACM porque es ahí donde el interés se coloca por primera vez en un entorno de formación institucionalizado. La pregunta pasa a ser entonces: ¿qué tipo de entorno educativo protege mejor el criterio de investigación a largo plazo?

Al hablar de la clase ACM de la Universidad Jiao Tong de Shanghái, Xie Saining subraya en particular la flexibilidad y la formación general. Por ejemplo, el «Foro de Estudiantes» (学子讲坛) pedía a los alumnos exponer cualquier tema ajeno a los cursos: podía ser filosofía, historia, sociedad o ciencia. El valor didáctico de este pasaje es el siguiente: si la formación temprana gira solo en torno a clasificaciones y a la optimización de problemas, es fácil obtener ejecutores eficaces; si permite a los estudiantes cultivar intereses amplios, es más probable que se forme un research taste.

También dice con claridad que no le gusta la competencia excesiva. No se opone a la competencia, sino a reducir a todos a un único indicador. El valor a largo plazo de la investigación suele venir de trayectorias no lineales: primero el interés y después el problema; primero experiencias a las que se llega por casualidad y después un hilo conductor que las une. La competencia excesiva acorta la escala temporal y lleva a perseguir solo las tareas que hoy otorgan puntos.

\begin{warningbox}{No hay que entender «persona común» como metas bajas}
The normal one no significa rebajar las metas, sino rechazar un relato mítico de uno mismo. Desplaza la atención de «¿soy o no un elegido?» a «¿puedo elegir de forma sostenida problemas que valgan la pena, encontrar personas con quienes valga la pena trabajar y asumir la incertidumbre a largo plazo?».
\end{warningbox}

\subsection{Resumen de la sección}

La historia de formación del inicio marca el tono de todo el episodio: a Xie Saining no le interesa ganar o perder en un punto aislado, sino la trayectoria a largo plazo. Esa trayectoria se manifestó después así: no seguir el camino más seguro de universidades prestigiosas y grandes empresas tecnológicas, y elegir una y otra vez a las personas y organizaciones más cercanas a sus propias preguntas de investigación.
\section{De Vision a Representation: el hilo conductor de la trayectoria de investigación}

La pista personal más importante de la entrevista es «el mundo nunca me deja hacer Vision, pero yo insisto en hacerla». Detrás de esta frase hay dos niveles de significado. El primero es la resistencia en las decisiones académicas y profesionales: al solicitar el doctorado casi no obtuvo ninguna de las vision offer que deseaba, y después logró entrar a UCSD gracias a un correo que le envió a Tu Zhuowen. El segundo es la resistencia del relato de la época: cuando los LLM se convirtieron en el centro, la computer vision pareció replegarse a la periferia. Sin embargo, el juicio de Xie Saining (谢赛宁) es que, si se entiende la vision como una perspective y no como un conjunto de tareas, entonces es precisamente el núcleo del camino hacia la inteligencia real.

\begin{figure}[H]
\centering
\includegraphics[width=0.96\textwidth]{trajectory-timeline.png}
\caption{La trayectoria de investigación en la entrevista con Xie Saining: de SJTU, UCSD, FAIR y DeepMind/NYU a AMI Labs.}
\end{figure}

\begin{knowledgebox}{Lectura de la figura: cada decisión filtra las condiciones de investigación}
Esta línea de tiempo no es un currículum, sino un mapa de la lógica de sus decisiones. La etapa en SJTU le dio una formación general y la entrada a la computación; la etapa en UCSD enfocó el problema en vision y representation; la etapa en FAIR, con experiencias como He Kaiming y ResNeXt, le permitió aprender el criterio de investigación de primer nivel; las etapas en DeepMind y NYU le mostraron la tensión entre el video, las tareas corporizadas y los recursos organizacionales; AMI Labs, por su parte, intenta impulsar el problema de los modelos del mundo dentro de una nueva forma de organización.
\end{knowledgebox}

\subsection{Cinco prácticas profesionales e investigación no lineal}

Durante el doctorado hizo prácticas en NEC Lab, Adobe, Meta, Google Research y DeepMind, entre otras. Algunas prácticas produjeron artículos; otras no dieron resultados evidentes. En la entrevista no presenta la «falta de productos» como un fracaso, sino como exploración: ir a distintas organizaciones para ver distintos problemas, entender distintas culturas de investigación y confirmar lo que no quería hacer.

Esto es especialmente importante para los investigadores jóvenes. Una línea de investigación no consiste en optimizar cada paso en función del número de artículos, sino en dejar que varias experiencias aparentemente dispersas acaben formando un problema interno. Si Xie Saining pudo después conectar vision, video, representation, world model y robotics, es precisamente porque en sus inicios no se limitó a optimizar de forma incremental dentro de un solo camino estrecho.

\begin{importantbox}{La investigación no es una point estimate, sino una integral en el tiempo}
En la entrevista usa la frase «no hay que preocuparse por la estimación en cada punto» para describir la evaluación de la investigación. Que un solo artículo, un solo rechazo o una sola práctica salgan bien o mal son apenas puntos en el eje del tiempo; la calidad real de un investigador se mide por la integral tras una acumulación de largo plazo.
\end{importantbox}

\subsection{FAIR, He Kaiming y ResNeXt}

La experiencia en FAIR es un punto decisivo de toda la entrevista. Xie Saining cuenta que, después de que He Kaiming se incorporó a FAIR, él colaboró con He en el ImageNet challenge durante el último mes de su práctica, y de ahí surgió el trabajo relacionado con ResNeXt. Lo importante aquí no es la historia de «encontrar a un mentor providencial», sino cómo un investigador de primer nivel pule una idea común hasta convertirla en una propuesta de representation escalable.

La idea de ResNeXt puede simplificarse así: sobre la base de ResNet se introduce la cardinality, que extiende una rama en varios group paralelos y obtiene una mejor capacidad de representación con un costo de cómputo similar. En la entrevista, Xie Saining también la relaciona con la intuición actual de MoE: la dispersión, el agrupamiento y la capacidad de escalar no son ideas que hayan surgido apenas hoy.

\begin{knowledgebox}{Términos clave: ResNeXt, cardinality y la intuición de MoE}
\begin{tabularx}{\textwidth}{p{0.22\textwidth}p{0.32\textwidth}X}
\toprule
Término & Problema que resuelve & Relación con el hilo conductor de este episodio \\
\midrule
ResNet & Facilita el entrenamiento de redes profundas mediante conexiones residuales & Representa el paradigma central del aprendizaje temprano de representaciones visuales. \\
ResNeXt & Aumenta la capacidad expresiva de la red con varias ramas paralelas agrupadas & Encarna «una representation más escalable con el mismo costo de cómputo». \\
Cardinality & Número de ramas o group paralelos & Está más cerca de un escalamiento estructurado que el simple aumento de profundidad o anchura. \\
MoE & Mixture of Experts, activación dispersa de distintos expertos & Comparte la intuición de ingeniería del agrupamiento y el escalamiento disperso. \\
\bottomrule
\end{tabularx}
\end{knowledgebox}

\subsection{Resumen de la sección}

La trayectoria de investigación de Xie Saining puede resumirse así: partió de tareas visuales concretas y poco a poco llegó a entender la vision como la suma de señales continuas de alta dimensión, representaciones jerárquicas, cognición espacial y temporal, e inteligencia predictiva. Esta definición prepara el terreno para los modelos del mundo que se tratan más adelante.

\section{Vision as a Perspective: la visión no es un campo menor}

La sección anterior vinculó la trayectoria de investigación personal con la representation; esta sección explica con claridad «por qué seguir trabajando en vision». Lo esencial aquí no es defender una disciplina antigua, sino redefinir el lugar de la vision en la siguiente generación de inteligencia.

Cuando la entrevista entra en «el mundo de las representaciones», Xie Saining (谢赛宁) plantea una distinción importante: la computer vision no es solo un conjunto de tareas como classification, detection y segmentation, sino una perspective desde la cual mirar la inteligencia. Trabaja con continuous, high-dimensional, noisy signals, es decir, señales en un espacio continuo, de alta dimensión y con ruido. Este tipo de señales es difícil de tokenizar de forma sencilla y tampoco trae de origen etiquetas escritas por personas.

Esto explica por qué el auge de los LLM no lo desanimó. El éxito de los LLM puso en primer plano las interfaces de lenguaje, los sistemas multimodales y el entrenamiento a mayor escala, y con ello le dio a la vision la oportunidad de dejar atrás las tareas aisladas y abordar los problemas amplios de la «inteligencia del mundo real». El verdadero peligro no es que los LLM sean demasiado potentes, sino que todos los problemas de visión se vean obligados a someterse a la narrativa de los modelos de lenguaje y que, al final, la visión quede reducida a un apéndice del prompt y del caption.

\begin{figure}[H]
\centering
\includegraphics[width=0.94\textwidth]{representation-ladder.png}
\caption{De las tareas de visión a los modelos predictivos del mundo: la escalera de capacidades que se desprende de la entrevista.}
\end{figure}

\begin{knowledgebox}{Lectura de la figura: la transición de capacidades de la task al world model}
Los niveles L0 a L4 de la figura no son una clasificación oficial, sino una reconstrucción de la lógica de la entrevista. Las primeras tareas de visión resuelven clasificación, detección y segmentación; la etapa multimodal convierte el lenguaje en la interfaz; más adelante, el sistema debe comprender flujos continuos de eventos, estructura espacial y las consecuencias de cada action. El punto de llegada no es «saber responder preguntas sobre una imagen», sino poder comprimir el flujo visual en una representación del mundo útil para predecir y planificar.
\end{knowledgebox}

\subsection{Por qué importan las representaciones jerárquicas}

Xie Saining menciona en repetidas ocasiones la hierarchical representation. La intuición es la siguiente: un agente no puede memorizar de forma explícita cada píxel, cada molécula y cada parámetro físico del mundo; tiene que aprender a abstraer. Abstraer no significa perder información, sino conservar la información relacionada con la tarea actual, con las acciones futuras y con el costo de las decisiones.

Por ejemplo, la mesa, el micrófono, la iluminación, el sonido y las texturas de una habitación pueden modelarse con todo detalle; pero si el objetivo es continuar la conversación, el sistema solo necesita conocer estados relevantes para la tarea: que el micrófono puede colocarse sobre la mesa, la posición relativa de las dos personas y si el sonido puede captarse. Este state «suficiente, no exhaustivo» es justamente el punto donde se encuentran el representation learning y el world model.

\begin{importantbox}{El problema central del aprendizaje de representaciones}
Una buena representation no reconstruye todos los detalles, sino que comprime señales de alta dimensión en un estado suficiente para sostener la predicción, la planificación y la acción. Debe estar más cerca del mundo físico que las etiquetas de lenguaje y, a la vez, ser más abstracta que los píxeles en bruto.
\end{importantbox}

\subsection{La ayuda y la contaminación del lenguaje}

El lenguaje como interface es sumamente útil. Permite que los sistemas multimodales definan problemas, formulen preguntas y den respuestas en lenguaje natural, y facilita que los sistemas de visión se alineen con los objetivos humanos. Pero el lenguaje también puede convertirse en un shortcut. En la entrevista, Xie Saining describe la tentación que el lenguaje ejerce sobre los sistemas de visión con las imágenes de una «muleta» y del «opio»: el lenguaje hace que el sistema parezca más inteligente, pero puede impedir que desarrolle la capacidad de procesar de verdad un mundo continuo.

\begin{warningbox}{Multimodal no equivale a comprensión real}
Si un benchmark de visión puede responderse correctamente apoyándose sobre todo en el sentido común lingüístico, entonces no demuestra que el modelo comprenda la imagen, el video o el mundo físico. Una vez que interviene el lenguaje, la evaluación debe verificar si el modelo realmente usó representaciones visuales o si resolvió la tarea solo con conocimiento previo textual.
\end{warningbox}

\subsection{Resumen de la sección}

La idea central de Vision as a perspective es esta: la visión no es una tarea menor absorbida por los LLM, sino el problema cognitivo de base que exige la construcción de modelos del mundo. Requiere que el modelo procese señales continuas, abstracciones jerárquicas, estructura espacial, flujos de eventos y las consecuencias reales de las acciones.
\section{Modelos del mundo: de la fórmula y el control al cerebro del robot}

La parte técnica más sólida de este episodio es la definición de modelo del mundo que da Xie Saining (谢赛宁). No trata world model como una palabra de moda; regresa a las preguntas fundamentales del control, del model-based reinforcement learning y de la ciencia cognitiva: dados el estado actual y una acción, predecir el estado siguiente, y después usar esa predicción para guiar la planning y la decision making.

\[
s_{t+1} = f(s_t, a_t)
\]

Aquí, $s_t$ representa el estado del sistema o del entorno en el instante actual; $a_t$ representa la action o intervention que realiza el agente; $f$ es la transition / predictive function aprendida, y $s_{t+1}$ es el estado siguiente al que el sistema puede llegar después de ejecutar la acción. La dificultad real no está en la fórmula, sino en cómo obtener un $s_t$ adecuado a partir de percepciones de alta dimensión, cómo lograr que $f$ aprenda las leyes físicas y cómo usar la predicción para actuar en lugar de solo generar muestras vistosas.

\begin{figure}[H]
\centering
\includegraphics[width=0.96\textwidth]{world-model-stack.png}
\caption{El ciclo cerrado mínimo de un modelo del mundo: percepción, representación del estado, predicción de la transición, planificación y decisión, retroalimentación real y actualización del aprendizaje.}
\end{figure}

\begin{knowledgebox}{Lectura de la figura: el modelo del mundo no es un módulo aislado}
Lo más importante de la figura es el ciclo cerrado. Primero, la entrada perceptual se comprime en un state relevante para la tarea; el modelo predice el estado siguiente a partir del state y la action; el planificador compara distintas trayectorias futuras; al ejecutar la acción se obtiene retroalimentación real, y los errores y fallas regresan al entrenamiento. Si falta cualquiera de estos eslabones, el world model se reduce a un generador estático, a un sistema de preguntas y respuestas o a un modelo de representación fuera de línea.
\end{knowledgebox}

\subsection{La intuición del Model Predictive Control}

La fórmula anterior muestra la forma mínima de un world model; ahora hay que responder cómo se convierte en acción. El MPC ofrece la respuesta más sencilla y, a la vez, la de mayor valor didáctico: usar el modelo para imaginar el futuro y ejecutar únicamente el paso que más conviene en el momento actual.

En la entrevista se menciona el model predictive control. Su procedimiento básico es el siguiente: en el instante actual, se usa el modelo para hacer el roll out de varias action sequence futuras, se calcula el cost de cada trayectoria, se elige la secuencia de menor cost, se ejecuta su primer paso y en el instante siguiente se vuelve a planificar. No se trata de «expresar razones» al estilo de un LLM, sino de elegir acciones con base en la predicción de estados futuros.

\begin{knowledgebox}{Términos clave: conceptos relacionados con los modelos del mundo}
\begin{tabularx}{\textwidth}{p{0.23\textwidth}p{0.30\textwidth}X}
\toprule
Término & Explicación breve & Función en este episodio \\
\midrule
State & Información mínima suficiente que describe el estado actual del sistema & Vincula el aprendizaje de representaciones con la predicción de acciones. \\
Action / Intervention & Acción o intervención que el agente ejerce sobre el entorno & Hace que la predicción pase de la observación a la decisión. \\
Transition Function & Función que predice el state siguiente a partir del state y la action actuales & Forma central del modelo del mundo. \\
Planning & Comparar trayectorias futuras y elegir una acción & Pone la predicción del modelo al servicio de un objetivo. \\
MPC & Model Predictive Control, predicción y control con horizonte deslizante & Muestra cómo se usa un modelo del mundo para el control. \\
Model-based RL & Uso de un modelo del entorno como apoyo del aprendizaje por refuerzo & Explica la relación histórica entre el world model y el RL. \\
\bottomrule
\end{tabularx}
\end{knowledgebox}

\subsection{El world model es un propósito, no un algoritmo único}

Xie Saining subraya que el modelo del mundo es difícil de definir porque no es el nombre de un algoritmo, sino un propósito. Los modelos de lenguaje, los video diffusion model, las 3D representation, la robotics, los VLA y el model-based RL pueden avanzar hacia el world model desde direcciones distintas. Discutir «cuál es el verdadero world model» tiene sentido a corto plazo, pero a largo plazo la pregunta más importante es si el sistema puede comprender el mundo físico, conservar la memory relevante, reason and plan, hacer counterfactual / causal inference y ser controllable and safe.

\begin{importantbox}{Las cinco capacidades de un modelo del mundo}
Un world model orientado al mundo físico debe tener, como mínimo: physical world understanding, large associated memory, reasoning / planning, counterfactual or causal inference, y controllability and safety. No se mide con un indicador tan simple como «qué tan realista se ve el video generado».
\end{importantbox}

\subsection{Resumen de la sección}

La definición de fondo del modelo del mundo es sencilla: predecir el state que sigue a una action. La dificultad está en la representación del state, en el anclaje físico de la predicción, en la utilidad práctica de la planning y en el ciclo cerrado de retroalimentación. Tratar el world model como un objetivo, y no como una arquitectura de modelo concreta, evita que las disputas terminológicas de corto plazo desvíen la atención.
\section{Contribuciones y límites de los LLM: el espacio virtual no es todo el mundo}

Ya se estableció la definición de modelo del mundo; ahora hay que distinguirlo de los LLM. Esta distinción es importante, porque la entrevista no está en contra de los LLM, sino que pregunta qué parte del mundo cubre realmente el éxito de los LLM.

Xie Saining (谢赛宁) no niega el carácter revolucionario de los LLM. Dice que tiene que agradecer a los LLM, porque sin ellos la inteligencia multimodal no habría escalado hasta su tamaño actual. Pero se opone a extender el relato de los LLM hasta afirmar que «los modelos de lenguaje conducen de manera natural a la human-level intelligence». En su marco, los LLM son más fuertes en el digital / virtual space: texto, código, síntesis de conocimiento, educación, derecho, búsqueda, agentic coding, etc. Pueden ser un elemento importante de un sistema inteligente, pero no son el fundamento del modelo del mundo.

\begin{figure}[H]
\centering
\includegraphics[width=0.96\textwidth]{llm-vs-world-model.png}
\caption{LLM y modelos del mundo: comparación de dos mercados de la inteligencia.}
\end{figure}

\begin{knowledgebox}{Lectura de la figura: los modelos de lenguaje y los modelos del mundo se complementan, pero no se sustituyen entre sí}
A la izquierda, los puntos fuertes del LLM son el conocimiento tokenizado, el razonamiento sobre texto y la operación en el espacio digital; a la derecha, el objetivo del world model son las señales continuas, la predicción física, la planificación de acciones y el control seguro. Ambos pueden complementarse: el lenguaje aporta la interfaz, el conocimiento y la comunicación; el modelo del mundo aporta el anclaje físico, la predicción dinámica y la capacidad de actuar. Reducir a la fuerza uno al otro lleva a diagnosticar mal los cuellos de botella tecnológicos.
\end{knowledgebox}

\subsection{Por qué los modelos de lenguaje se parecen al aprendizaje con supervisión fuerte}

En la entrevista aparece una idea muy reveladora: a los modelos de lenguaje se les suele describir como self-supervised learning, pero desde otro ángulo, el lenguaje mismo ya es una señal de supervisión fuerte, procesada por la civilización humana. Miles de años de civilización, libros, páginas web, código, artículos científicos e internet han comprimido una enorme cantidad de conocimiento del mundo en texto tokenizado. Entrenar un LLM se parece a descargar ese conocimiento ya procesado, no a aprender directamente del mundo físico.

Esto explica por qué la scaling law de los LLM pudo aparecer relativamente pronto: cuenta con una gran cantidad de material que los seres humanos ya organizaron, etiquetaron y volvieron comunicable. En cambio, vision y robotics trabajan con raw sensory data, acciones, espacio, restricciones físicas y retroalimentación, y por naturaleza carecen de «etiquetas al estilo de internet» de la misma escala, la misma calidad y el mismo grado de compresión.

\begin{warningbox}{«Datos gratuitos» no equivale a «datos sin supervisión»}
Que el texto de internet pueda recopilarse gratis no significa que carezca de etiquetas humanas. El lenguaje es una estructura que los seres humanos comprimieron durante largo tiempo para comunicarse; cuando el modelo ingiere ese texto, también ingiere la gran cantidad de abstracción, selección y etiquetado que los seres humanos ya realizaron.
\end{warningbox}

\subsection{Una taza que se rompe: lo que omite la descripción lingüística}

Xie Saining da un ejemplo sencillo: cuando decimos «la taza cayó al suelo y se rompió», el lenguaje solo conserva la información necesaria para comunicarse. No describe el contacto de la taza, las fuerzas que actúan sobre ella, la trayectoria de la fractura, las propiedades del material, la dispersión de los fragmentos ni el sonido. Para la comunicación humana, la mayoría de esos detalles son innecesarios; para un sistema que debe actuar en el mundo físico, algunos de ellos pueden ser cruciales.

Por lo tanto, el lenguaje no es el mundo mismo, sino una compresión hecha para comunicarse. Lo que el modelo del mundo necesita aprender son precisamente las regularidades que quedan fuera de esa compresión lingüística: dinámica, estructura espacial, contacto, secuencia temporal, causalidad y posibilidad de acción.

\subsection{Resumen de la sección}

El éxito de los LLM proviene de la enorme compresión del conocimiento humano y de la escala de internet. Seguirán siendo un componente importante de los sistemas inteligentes, pero la inteligencia orientada al mundo físico no puede depender solo del lenguaje. Lo que el modelo del mundo debe completar es la parte del mundo que el lenguaje omite por naturaleza.
\section{De «descargar internet» a «descargar a los humanos / el mundo»}

Una de las metáforas que más se difundieron de la entrevista es la del «OpenAI inverso». La ruta directa es: descargar datos de internet, entrenar un transformer / GPT, obtener inteligencia lingüística y después llevarla al mercado. Esta ruta funciona porque internet ya contiene una gran cantidad de conocimiento escrito por personas. Los modelos del mundo no cuentan con el mismo shortcut: los datos del mundo real, los escenarios empresariales, los procesos físicos, las tareas robóticas y la retroalimentación de sensores no pueden simplemente descargarse de páginas web.

\begin{figure}[H]
\centering
\includegraphics[width=0.96\textwidth]{download-internet-to-world.png}
\caption{De «descargar internet» a «descargar a los humanos / el mundo»: los modelos del mundo necesitan entornos reales y un ciclo cerrado con socios.}
\end{figure}

\begin{knowledgebox}{Lectura de la figura: lo esencial de la ruta inversa es el ciclo cerrado de colaboración}
A la izquierda está la ruta directa, de internet a los modelos de lenguaje; a la derecha, la ruta inversa, en la que el mundo, los socios, las tareas y los datos retroalimentan en conjunto al modelo. Las condiciones del ciclo cerrado en la parte inferior de la figura indican que un modelo del mundo necesita datos de socios, un modelo inicial, tareas reales, retroalimentación de evaluación y control de seguridad. Sin la participación del mundo, el modelo del mundo se queda en un concepto de artículo académico o en una demo generativa.
\end{knowledgebox}

\subsection{World model needs the world}

Xie Saining (谢赛宁) dice: world model needs the world. Esta frase es el núcleo de la estrategia organizacional de AMI Labs. Significa que el modelo del mundo no es una tarea que una sola empresa pueda completar con datos cerrados, sino que requiere la colaboración de granjas, hospitales, fábricas, empresas de robótica, redes de sensores, entornos de simulación y escenarios de negocio reales.

En este marco, el modelo aporta primero una capacidad inicial y entra en escenarios reales para generar valor; los escenarios reales producen retroalimentación, errores y datos nuevos; y esos datos retroalimentan a su vez al modelo. Este ciclo cerrado se parece a un motor de datos, pero su objeto ya no es solo el texto de páginas web, sino el mundo físico y los sistemas de acción.

\begin{importantbox}{La esencia del OpenAI inverso}
El OpenAI directo descarga internet; el OpenAI inverso tiene que construir los datos junto con el mundo. El primero depende de corpus de texto ya existentes; el segundo depende de tareas reales, redes de socios, sistemas de evaluación y retroalimentación continua.
\end{importantbox}

\subsection{Por qué esto no es simplemente comprar datos}

Decir que world model needs the world puede hacer pensar que la solución consiste solo en adquirir más datos físicos. Esta sección desmonta ese malentendido: lo que de verdad necesita un modelo del mundo es un mecanismo de generación de datos, no un inventario de datos de una sola vez.

Si lo que necesita el modelo del mundo son las consecuencias de cada action, los estados físicos, la recuperación ante fallas, los flujos de sensores y los procesos de cada dominio, entonces «comprar un lote de datos» está muy lejos de ser suficiente. Los datos deben tener tareas, objetivos, retroalimentación, evaluación y actualización continua. Más importante aún, muchos datos solo aparecen después de que el modelo entra en el escenario: el modelo se equivoca, una persona lo corrige, el sistema se recupera, el usuario cambia el flujo de trabajo; nada de esto puede prepararse de antemano en un corpus estático.

\begin{knowledgebox}{De data factory a environment factory}
Este episodio puede leerse junto con el panorama de datos de EP134: lo que necesita un modelo del mundo no son datos adquiridos en una sola compra, sino entornos capaces de producir de forma continua tareas, fallas, retroalimentación y evaluaciones. La fábrica de datos del futuro se parecerá más a una environment factory.
\end{knowledgebox}

\subsection{Resumen de la sección}

«Descargar internet» resolvió la primera etapa de la inteligencia lingüística; «descargar a los humanos / el mundo» implica pasar del texto estático a los entornos dinámicos. La competencia en modelos del mundo no es solo una competencia de parámetros del modelo, sino también de redes de colaboración, sistemas de evaluación y ciclos cerrados de retroalimentación real.
\section{Research taste y el \emph{Sutra del Diamante}: cómo juzgar si un problema vale la pena}

A la mitad de la entrevista hay un tramo largo dedicado a la research taste, en el que se menciona el verso del \emph{Sutra del Diamante} (《金刚经》) «como un sueño, una ilusión, una burbuja, una sombra». Esta parte no es un paréntesis místico, sino metodología de investigación. Lo que le interesa a Xie Saining (谢赛宁) es cómo mantener el juicio sobre el problema mismo en un entorno de investigación de gran incertidumbre, con una evaluación muy ruidosa y resultados de corto plazo que con frecuencia engañan.

\begin{figure}[H]
\centering
\includegraphics[width=0.92\textwidth]{research-taste-map.png}
\caption{Los componentes de la research taste: elección del problema, criterio estético, evidencia y momento oportuno.}
\end{figure}

\begin{knowledgebox}{Lectura de la figura: la research taste no es inspiración, sino un sistema de juicio a largo plazo}
Las cuatro esquinas de la figura son la elección del problema, el criterio estético, la evidencia y el momento oportuno. Para juzgar qué problema vale la pena, un investigador no puede fijarse solo en lo que está de moda ni en si puede publicar un artículo rápido; también debe preguntarse si el problema seguirá teniendo sentido dentro de diez años, si la evidencia respalda la dirección, si los recursos de la organización son suficientes y si él mismo puede tolerar una incertidumbre prolongada.
\end{knowledgebox}

\subsection{Criterio estético: contra las citas baratas y los eslóganes}

El párrafo anterior descompuso la research taste en problema, evidencia, criterio estético y momento oportuno; aquí se desarrolla primero el «criterio estético». No se trata de una preferencia de estilo, sino de la sensibilidad del investigador ante el abuso de conceptos, el empaque con eslóganes y la sustitución encubierta de la evidencia.

Al final, Xie Saining critica que se tome la frase de Wittgenstein (维特根斯坦) «los límites de mi lenguaje son los límites de mi mundo» para avalar directamente a los LLM, y también que se use a la ligera la frase de Feynman (费曼) «What I cannot create, I do not understand» para avalar los unified model. No se opone a la filosofía ni a las citas célebres, sino a las citas decorativas sacadas de contexto.

Detrás de esto está el criterio estético de la research taste: un argumento no puede sostenerse en citas de personajes célebres, sino en definiciones, supuestos, mecanismos y evidencia. La filosofía puede inspirar la investigación, pero no sustituye a un technical argument.

\begin{warningbox}{Una cita célebre no es un argumento técnico}
Poner una frase célebre de filosofía o de física al inicio de un artículo no demuestra por sí solo que la línea de modelos elegida sea correcta. Una investigación de calidad necesita explicar qué significa el concepto en su contexto original, si los supuestos siguen siendo válidos al trasladarlo a la IA y qué mecanismos y evidencia respaldan ese traslado.
\end{warningbox}

\subsection{Visión de largo plazo: del rechazo a la prueba del tiempo}

En la entrevista también se cuenta que un artículo fue rechazado en su momento por cuestiones de detalle, y que después se publicó en otro congreso y obtuvo un test-of-time award. La historia muestra que el valor de una investigación y el resultado de la revisión no siempre van de la mano. La revisión de corto plazo puede fijarse en exceso en el formato, los detalles o las preferencias dominantes del momento; el impacto de largo plazo depende más de que el problema sea real, de que el método sea reutilizable y de que la idea se incorpore a trabajos posteriores.

\begin{importantbox}{La escala de tiempo de la evaluación de la investigación}
La evaluación de corto plazo mira el accept / reject; la de largo plazo mira si la idea cambia la definición del problema, las herramientas metodológicas o el lenguaje de la comunidad. La verdadera research taste consiste en mantener, en medio del ruido de corto plazo, una dirección reutilizable a largo plazo.
\end{importantbox}

\subsection{Resumen de la sección}

La research taste es lo que une la línea técnica y la línea personal de este episodio. Explica por qué Xie Saining eligió vision, eligió FAIR, dejó un camino seguro y fundó una empresa para trabajar en world model, y explica también por qué le disgusta dejarse llevar por eslóganes y palabras de moda.
\section{AMI Labs: la forma de organización también es parte de la ruta técnica}

Cuando la entrevista llega a AMI Labs, la pregunta deja de ser «qué es un modelo del mundo» y pasa a ser «qué organización puede construir un modelo del mundo». Xie Saining (谢赛宁) considera que en la universidad no alcanzan los recursos; que en las grandes empresas los ciclos de producto y la carrera armamentista por las tablas de clasificación reducen el espacio para explorar; que los institutos de investigación tradicionales tienen libertad, pero difícilmente pueden asumir entrenamientos a gran escala ni colaboraciones en escenarios reales, y que las empresas de grandes modelos completamente cerradas pueden cortar el vínculo con la academia y la discusión abierta. Por eso, AMI Labs intenta encontrar una forma intermedia.

\begin{figure}[H]
\centering
\includegraphics[width=0.96\textwidth]{ami-labs-thesis.png}
\caption{La tesis organizacional de AMI Labs: buscar el equilibrio entre la línea principal del modelo del mundo, una organización favorable a la investigación, una red global de socios, un ciclo comercial cerrado y el vínculo con la academia.}
\end{figure}

\begin{knowledgebox}{Lectura de la figura: cómo el diseño organizacional influye en las posibilidades técnicas}
Los cinco nodos de la figura forman en conjunto la tesis de AMI. El modelo del mundo requiere research de largo plazo; el research de largo plazo requiere recursos y oxígeno organizacional; los datos del mundo real requieren una red global de socios; un emprendimiento serio requiere un ciclo comercial cerrado, y la línea de investigación necesita mantener el vínculo con la academia. Si falta cualquiera de estos nodos, el world model puede degradarse en un proyecto de artículos, en un producto cerrado o en una visión vacía.
\end{knowledgebox}

\subsection{Ni instituto de investigación puro ni empresa de producto cerrada}

La figura anterior ya presentó los cinco nodos de AMI; esta sección examina primero la tensión organizacional central. El modelo del mundo necesita libertad de investigación, pero una empresa emergente tiene que lidiar con presiones de recursos, clientes, capital y entregas.

Xie Saining explica que AMI no es una non-profit ni un research lab puro; necesita un business model. Pero tampoco quiere convertirse en una empresa cerrada que solo gire en torno a los deadline de producto y a los benchmark. Esta tensión es difícil, porque una empresa comercial tiene que responder por sus recursos y sus resultados, mientras que el modelo del mundo es una dirección altamente exploratoria.

Por eso él insiste en una «researcher friendly organization». Si la organización no tiene suficiente oxígeno, aunque los investigadores sepan que cierto problema es importante, terminarán asignados a eslabones entregables de la cadena de producto de corto plazo, como video captioning, optimización para tablas de clasificación o soporte a los ciclos de lanzamiento.

\begin{importantbox}{La forma de organización cambia el espacio de investigación}
El mismo investigador, con la misma idea, puede hacer cosas completamente distintas en una universidad, en una gran empresa, en una empresa emergente cerrada o en una empresa emergente favorable a la investigación. Un problema de la magnitud del modelo del mundo necesita recursos, libertad para explorar y, además, colaboración en escenarios reales.
\end{importantbox}

\subsection{El papel de Yann LeCun: ruta, carácter y confianza}

Una vez expuesta la tensión organizacional, falta ver quién la sostiene. En la entrevista, Yann LeCun no es solo el nombre de un cofundador, sino la combinación de convicción sobre la ruta, integrity de científico y estabilizador psicológico del equipo.

A lo largo de la entrevista, Xie Saining habla varias veces de Yann LeCun. Subraya que LeCun no se opone a los LLM, sino a la narrativa de que «los LLM conducen naturalmente a una inteligencia de nivel humano». Lo que admira de LeCun es, primero, su apego de largo plazo a la ruta de world model / JEPA / autonomous intelligence; segundo, su integrity como científico, y tercero, su temperamento personal: ama la vida, tiene un mundo amplio y logra que quienes lo rodean sientan que hay un camino por delante.

Este pasaje deja una lección para las organizaciones técnicas: un founder o un líder científico no solo aporta dirección, también aporta la psicología de la organización. Cuanto más incierto es el problema de investigación, más necesita el equipo a alguien que ofrezca convicción de largo plazo, corrección oportuna y un espacio de discusión donde se permita cuestionar.

\subsection{Resumen de la sección}

La historia de AMI Labs muestra que la ruta técnica y la ruta organizacional son inseparables. El modelo del mundo no avanza solo gracias a una arquitectura de modelo; necesita el sostén conjunto de recursos, socios, cultura de investigación, un ciclo comercial cerrado y una narrativa de largo plazo.
\section{«Silicon Valley is LLM-pilled»: narrativa, rankings y asignación de recursos}

Una vez expuestas las decisiones organizativas de AMI, se entiende por qué la entrevista dirige sus críticas contra la «narrativa de Silicon Valley». El objeto de la crítica no es un lugar geográfico, sino un sistema industrial que vincula entre sí los problemas, los rankings y los recursos.

El «Silicon Valley está hipnotizado» del título de la entrevista se refiere a que toda la industria de la IA está fuertemente organizada en torno a la narrativa de los LLM. Xie Saining (谢赛宁) no afirma que los LLM carezcan de valor, sino que una narrativa define el benchmark, el benchmark define la resource allocation y la resource allocation, a su vez, determina lo que los investigadores pueden hacer. Al final, muchas personas capaces no es que no quieran trabajar en modelos del mundo, comprensión de video o inteligencia física, sino que los objetivos de la organización las asignan a posiciones más cercanas a la cadena de valor actual.

\subsection{Cómo la cadena de valor reduce el espacio de investigación}

La cadena descrita en la entrevista puede resumirse así:

\begin{center}
\begin{tabularx}{0.95\textwidth}{p{0.18\textwidth}p{0.28\textwidth}X}
\toprule
Eslabón & Manifestación & Consecuencia \\
\midrule
Narrativa & AGI, scaling law, LLM frontier & Define qué problemas «parecen importantes». \\
Rankings & Chatbot Arena, clasificaciones de matemáticas, código y capacidades generales & Orienta a los equipos a optimizar métricas visibles. \\
Asignación de recursos & El cómputo, el personal y el ritmo de lanzamientos se concentran en los rankings y en los productos & La investigación exploratoria se queda sin oxígeno. \\
Asignación de puestos & Problemas como la comprensión de video se fragmentan en tareas de captioning o de soporte de producto & Resulta difícil desarrollar una verdadera ruta world-model-first. \\
\bottomrule
\end{tabularx}
\end{center}

\begin{warningbox}{El riesgo de quedar hipnotizado por una narrativa}
Cuando una narrativa es demasiado fuerte, los investigadores confunden «útil para el ranking» con «útil para la naturaleza de la inteligencia». No es un problema de una empresa en particular, sino el efecto de filtrado que toda la cadena de valor de la industria ejerce sobre los problemas de investigación.
\end{warningbox}

\subsection{Por qué escapar de Silicon Valley no es una cuestión geográfica}

«Escapar de Silicon Valley» no significa simplemente mudarse a otro lugar. Silicon Valley sigue concentrando la mayor densidad de talento, capital y cultura de ingeniería, y es posible que AMI abra en el futuro una sede ahí. De lo que realmente hay que escapar es del control que ejercen sobre la definición de los problemas la narrativa única de los LLM, los ciclos de lanzamiento de productos y la carrera armamentista de los rankings.

En la entrevista, las oficinas en varias ciudades (París, Nueva York, Montreal y Singapur) no son solo una disposición administrativa: también responden a la idea organizativa de que «world model needs the world», es decir, a la participación conjunta de distintas regiones, industrias, redes académicas y fuentes de datos.

\subsection{Resumen de la sección}

El problema central de estar LLM-pilled no es que los LLM sean demasiado potentes, sino que la narrativa es demasiado única. La ruta de los modelos del mundo necesita abrirse espacio junto a la cadena de valor existente, para que el video, la física, la robótica, los procesos industriales y la retroalimentación real vuelvan a ser objetos de primer orden en la investigación en IA.
\section{Robots, VLA y la segunda mitad del preentrenamiento}

En la segunda parte de la entrevista se vuelve varias veces a la robotics. El juicio de Xie Saining (谢赛宁) es que, antes de hablar de AGI o de super intelligence, primero hay que preguntarse si se puede construir un robot lo bastante confiable para encargarse de las tareas domésticas en un hogar. Muchas cosas que puede hacer un niño de pocos años o un adolescente, los robots de hoy todavía no las hacen bien. No se trata de un problema aislado del hardware de las extremidades, sino de un problema del cerebro del robot.

\subsection{La división del trabajo entre VLA y los modelos del mundo}

VLA significa Vision-Language-Action; su objetivo es conectar la visión, las instrucciones en lenguaje y las acciones. Es una forma importante dentro de la ruta de la robótica, pero Xie Saining considera que no basta con tomar un language model como foundation y añadirle una action head para resolver el problema del world model pre-training. Un VLA puede ser muy fuerte en tareas concretas y, aun así, no necesariamente asumir el trabajo de la capa base de la «segunda mitad del preentrenamiento».

\begin{knowledgebox}{Términos clave: algunas palabras clave en la ruta de la robótica}
\begin{tabularx}{\textwidth}{p{0.22\textwidth}p{0.32\textwidth}X}
\toprule
Término & Problema que resuelve & Significado en este episodio \\
\midrule
VLA & Conexión de extremo a extremo entre visión, lenguaje y acción & Una ruta importante para que los robots ejecuten tareas, pero no equivale a un modelo del mundo completo. \\
Robot Brain & El cerebro del robot o su capa de inteligencia superior & Requiere percepción, memoria, predicción, planificación y control. \\
Hardware Scaling Law & Obtener más datos reales y más experiencia de hardware mediante el despliegue de más robots & Las empresas que fabrican el cuerpo del robot deben enfrentarla, pero no resuelve directamente el preentrenamiento del cerebro. \\
Imitation Learning & Aprender políticas de acción a partir de trayectorias de demostración & Útil a corto plazo, pero depende de los datos y de la distribución de tareas. \\
World Model Pre-training & Preentrenamiento base orientado a señales multimodales continuas & Lo que Xie Saining llama la segunda mitad del preentrenamiento. \\
\bottomrule
\end{tabularx}
\end{knowledgebox}

\subsection{Qué entra y qué sale en la segunda mitad del preentrenamiento}

El conductor insiste en preguntar qué entra y qué sale en el world model pre-training. La respuesta de Xie Saining es que, al menos a largo plazo, la entrada debería ser una señal multimodal de espacio continuo, de alta dimensión y con ruido; al principio puede ser video, y tal vez incluya otros encoder además del visual. Qué debe ser la salida sigue siendo una research question. Este «no lo sé» es importante: indica que el problema todavía no se ha reducido prematuramente a una recipe fija.

\begin{importantbox}{El carácter abierto de la segunda mitad del preentrenamiento}
El preentrenamiento de los modelos de lenguaje tiene una forma clara: next-token prediction. El preentrenamiento de los modelos del mundo todavía no cuenta con un paradigma igual de estable. Las entradas posibles son video, sensores y flujos multimodales; los objetivos posibles son la predicción, la representación, la sorpresa, las consecuencias de las acciones o estados sobre los que se pueda planificar. Aquí siguen abiertas preguntas de investigación básica.
\end{importantbox}

\subsection{Resumen de la sección}

Los robots son una de las salidas más naturales de los modelos del mundo, pero los recursos de corto plazo de las empresas de robótica suelen estar dirigidos por el despliegue de hardware y por tareas concretas. Lo que la ruta de los modelos del mundo debe aportar es el cerebro de la capa superior: capaz de entender el mundo físico, conservar memoria, predecir consecuencias y servir de base para los VLA y el control de robots.
\section{La inteligencia no es una línea de puntaje en lenguaje}

Al final de la entrevista, Xie Saining (谢赛宁) sostiene que la AGI es un falso problema y habla de la inteligencia animal y de la arrogancia humana. Cita libros sobre cognición animal y señala que cada ser vivo tiene sus propias formas de percibir y de actuar: el razonamiento de los chimpancés, el almacenamiento de alimento de las aves, la comunicación de las ballenas, el olfato de los perros, el oído de los murciélagos. La implicación técnica es que la inteligencia no se ordena sobre una línea de puntaje de un examen de lenguaje, sino que está ligada al cuerpo, al entorno, a los objetivos, a los canales de percepción y a las tareas de supervivencia.

\begin{figure}[H]
\centering
\includegraphics[width=0.94\textwidth]{intelligence-spectrum.png}
\caption{La inteligencia no es una línea de puntaje en lenguaje: del texto y el código a la visión, la acción, el cuerpo y el entorno, y los objetivos sociales.}
\end{figure}

\begin{knowledgebox}{Lectura de la figura: por qué importa una perspectiva no antropocéntrica}
La figura descompone la inteligencia en varias dimensiones: texto/código, visión/espacio, acción/control, cuerpo/entorno y sociedad/objetivos. Que un LLM rinda bien en texto y código no significa que posea automáticamente capacidades en el plano del cuerpo, del entorno y de la acción. El sentido de los modelos del mundo es precisamente devolver la inteligencia a la percepción, la predicción, la acción y la retroalimentación.
\end{knowledgebox}

\subsection{La inteligencia de una ardilla y las tareas domésticas de un niño}

En la entrevista se cita una perspectiva de Rich Sutton: en lugar de maravillarse porque un modelo escribe código o gana medallas en competencias de matemáticas, conviene pensar en lo difícil que es construir una ardilla capaz de sobrevivir en el mundo real. La afirmación parece exagerada, pero desplaza la pregunta de la «capacidad para aprobar exámenes humanos» a la capacidad básica de «un ser vivo que actúa en el mundo».

Xie Saining añade que un niño de pocos años, o uno de doce, puede hacer muchas tareas domésticas que los robots actuales todavía no realizan de forma confiable. No se trata de menospreciar a los LLM, sino de recordar que la percepción, el movimiento, la recuperación ante errores, los objetivos y el sentido común en el mundo real siguen siendo grandes problemas sin resolver.

\begin{warningbox}{El lema de la AGI tiende a ocultar las capacidades concretas}
Si no se definen el entorno, el cuerpo, los objetivos y la evaluación, la AGI se convierte con facilidad en una etiqueta vacía. Una pregunta más operable es: ¿puede el sistema, en un entorno real determinado, percibir, predecir y actuar de forma confiable, recuperarse después de equivocarse y transferir gradualmente lo aprendido a tareas nuevas?
\end{warningbox}

\subsection{Resumen de la sección}

La inteligencia no tiene una sola escala. La capacidad lingüística es importante, pero la inteligencia en el mundo real también abarca la visión, el cuerpo, la acción, la retroalimentación y los objetivos sociales. La línea de investigación de los modelos del mundo busca cubrir justamente estas dimensiones que el puntaje en lenguaje deja ocultas.
\section{«42»: citas filosóficas, destino y los límites del modelo del mundo}

El «42» del final viene de una broma de «Guía del viajero intergaláctico»: la respuesta a la pregunta fundamental de la vida, el universo y todo lo demás podría ser 42. Al cierre de la entrevista, el conductor pregunta «¿este mundo es un enorme modelo del mundo?» y «¿puedes predecir el destino?». Xie Saining (谢赛宁) responde que, por supuesto, el mundo puede verse como un enorme modelo del mundo, pero que no podemos predecir el destino porque no alcanzan los recursos; tal vez habría que usar la Tierra o el universo entero como computadora para obtener la respuesta.

Este pasaje no es un cierre en broma, sino un regreso a los límites del modelo del mundo: un world model existe para hacer predicciones suficientemente buenas con recursos limitados, no para copiar el universo de forma omnisciente y omnipotente. Cuanto más se acerca el modelo al mundo real, más tiene que enfrentar las limitaciones de recursos de cómputo, compresión de la representación, elección de objetivos e incognoscibilidad.

\subsection{Soltar a Wittgenstein}

Antes, el «42» sirvió para recordar que el modelo del mundo no puede fingir omnisciencia; aquí se vuelve a la relación entre el lenguaje y el mundo. Esta subsección importa porque muchos relatos sobre los LLM recurren justamente a frases filosóficas célebres para presentar la capacidad lingüística como comprensión del mundo.

La crítica de Xie Saining al uso de las citas de Wittgenstein es, en realidad, una forma de trazar los límites del modelo del mundo. La tesis del primer Wittgenstein sobre lenguaje y mundo tiene un contexto filosófico específico, y el Wittgenstein tardío se orientó después hacia los juegos de lenguaje. Tomar «los límites de mi lenguaje son los límites de mi mundo» como prueba directa de que los LLM pueden abarcar el mundo es un desajuste de contexto.

Si se sigue la idea tardía de los «juegos de lenguaje», según la cual el significado del lenguaje proviene de la práctica y de las formas de vida, esto se acerca más bien a la postura del modelo del mundo: el lenguaje no son tokens suspendidos en el aire, sino que adquiere significado al actuar, usarse y recibir retroalimentación en el mundo real.

\begin{importantbox}{La relación entre el lenguaje y el mundo}
El lenguaje puede comprimir el mundo, comunicarlo y organizar la experiencia, pero no puede sustituirlo. Lo que el modelo del mundo debe aprender es la parte de la estructura en la que el lenguaje se relaciona con la práctica: acción, retroalimentación, restricciones, causalidad y previsibilidad.
\end{importantbox}

\subsection{The normal one y la batería del equipo}

Por último, se vuelve a the normal one. Xie Saining retoma una expresión de Jürgen Klopp y se imagina como la batería del equipo: alguien que, con entusiasmo y energía, les da carga a los demás. No se trata de un discurso motivacional vacío, sino de una función central de una research organization. La exploración a largo plazo trae mucha frustración, fracasos e incertidumbre, y el equipo necesita a alguien que sostenga la dirección, la confianza y el ritmo.

También reconoce que el trasfondo de la investigación suele ser desalentador: los momentos de verdadera alegría quizá sean solo del 5\% al 10\%, cuando algo por fin sale. La afirmación es muy honesta y también explica por qué una línea como la del modelo del mundo necesita el respaldo de una cultura organizacional. Sin resiliencia psicológica a largo plazo, una visión ambiciosa pronto queda aplastada por los reveses de corto plazo.

\subsection{Resumen de la sección}

El cierre junta las preguntas técnicas y las preguntas de la vida: el mundo puede modelarse, pero el modelo siempre está limitado por los recursos y los objetivos; el lenguaje puede inspirar el pensamiento, pero no sustituye la práctica; una persona común también puede llegar a los grandes problemas mediante decisiones de largo plazo, la energía del equipo y el criterio de investigación.
\section{Términos clave: índice de palabras clave de este episodio}

\begin{longtable}{p{0.24\textwidth}p{0.32\textwidth}p{0.35\textwidth}}
\toprule
Término & Explicación en una frase & Papel en este episodio \\
\midrule
World Model & Modelo que predice estados futuros a partir del estado y la acción & Tesis técnica central de la entrevista. \\
Predictive Brain & Cerebro predictivo orientado al mundo físico & Capa de inteligencia de nivel superior que AMI Labs quiere construir. \\
State & Representación comprimida del entorno, relevante para la tarea & Vincula el aprendizaje de representaciones con el planning. \\
Action & Acción o intervención que el agente ejerce sobre el entorno & Hace que la predicción intervenga en el control y la toma de decisiones. \\
Representation Learning & Aprendizaje de representaciones abstractas reutilizables & Puente que lleva de la vision al world model. \\
Vision as Perspective & Entender la visión como la perspectiva fundamental de la inteligencia, no como un conjunto de tareas & Explica el sentido de la CV en la era de los LLM. \\
LLM-pilled & Que la narrativa de los LLM organiza en exceso los recursos y la imaginación & Crítica a la cadena de valor de Silicon Valley. \\
VLA & Modelo Vision-Language-Action & Forma importante, pero insuficiente, de la línea de la robótica. \\
World Model Pre-training & Preentrenamiento fundacional orientado a video, sensores y el mundo físico & Problema abierto central de la segunda mitad del preentrenamiento. \\
Research Taste & Capacidad de juzgar problemas, evidencia, momento oportuno y estética & Explica las decisiones de personas y organizaciones. \\
JEPA & Joint Embedding Predictive Architecture & Idea representativa de la línea de LeCun: «predecir en un espacio de representaciones abstractas». \\
OpenAI a la inversa & Construir modelos junto con el mundo y una red de socios, en lugar de solo descargar internet & Metáfora organizacional y comercial de AMI. \\
\bottomrule
\end{longtable}

\subsection{Resumen de la sección}

Todas estas palabras clave apuntan a la misma pregunta: si en la siguiente etapa la IA podrá pasar de «saber procesar el mundo que los humanos escribieron» a «poder predecir, actuar y aprender en el mundo real». Esto exige actualizar al mismo tiempo los modelos, los datos, la organización y los límites filosóficos.

\section{Síntesis y ampliación}

\subsection{Conclusiones centrales}

\begin{enumerate}
\item La trayectoria personal de Xie Saining (谢赛宁) no es un manual de éxito, sino una muestra de cómo un problema de investigación toma forma de manera gradual mediante elecciones, organización y azar.
\item La visión no debe entenderse como un conjunto de tareas antiguas, sino como una perspective para procesar señales del mundo real continuas, de alta dimensión y con ruido.
\item La definición mínima de un modelo del mundo es $s_{t+1}=f(s_t,a_t)$, pero lo verdaderamente difícil es la representación del state, el anclaje físico, el planning y el ciclo cerrado de retroalimentación.
\item El LLM es un componente importante de un sistema inteligente, pero se alimenta sobre todo del conocimiento que la humanidad ya ha tokenizado y no puede sustituir de forma automática el modelado del mundo físico.
\item «OpenAI a la inversa» significa que un modelo del mundo no puede limitarse a descargar internet, sino que debe construir junto con socios del mundo real ciclos cerrados de datos, evaluación y aplicación.
\item Lo esencial de AMI Labs no es ser «otra empresa de modelos grandes», sino intentar una nueva forma de organización que equilibre libertad de investigación, escala de recursos, ciclo comercial cerrado y vínculo con la academia.
\item El problema de estar LLM-pilled es que una narrativa única define las tablas de clasificación y la asignación de recursos, con lo que reduce el espacio de exploración de direcciones como los modelos del mundo, la comprensión de video y la inteligencia física.
\item La inteligencia no es una sola línea de puntuación lingüística; la inteligencia del mundo real tiene que volver al cuerpo, el entorno, la acción, la retroalimentación y los objetivos.
\end{enumerate}

\subsection{Preguntas abiertas}

Después de la síntesis conviene mantener una actitud de interrogación, porque uno de los aspectos más honestos de este episodio es reconocer que muchas respuestas clave sobre los modelos del mundo todavía no han convergido. Las siguientes preguntas son también pistas que deben seguirse de manera continua en lecturas posteriores y al generar nuevas notas.

\begin{itemize}
\item ¿El preentrenamiento de modelos del mundo dará lugar a un paradigma estable similar a la next-token prediction?
\item ¿El state de un modelo del mundo debe ser 3D explícito, una latent representation, una estructura híbrida o formarse de manera dinámica según la tarea?
\item ¿La frontera entre VLA y world model se irá fusionando o se formará una división del trabajo entre una capa base aguas arriba y una capa de acción aguas abajo?
\item ¿Una red global de socios al estilo de AMI podrá aportar realmente un ciclo cerrado de datos insustituible?
\item Cuando las tablas de clasificación de LLM se saturen, ¿la industria volverá a dirigir recursos hacia la physical intelligence?
\end{itemize}

\subsection{Lecturas adicionales}

\begin{itemize}
\item Yann LeCun, \textit{A Path Towards Autonomous Machine Intelligence}: para entender la ruta de largo plazo de JEPA, los modelos del mundo y la autonomous intelligence.
\item Richard Sutton, artículos sobre Dyna / model-based reinforcement learning: para entender la relación histórica entre los modelos del mundo, el planning y el RL.
\item Materiales introductorios de Model Predictive Control: para entender la intuición de control de «usar un modelo para hacer roll out del futuro y elegir una acción».
\item Artículos sobre aprendizaje de representaciones visuales, generación de video, 3D representation, VLA y preentrenamiento de robots: para entender las múltiples vías de entrada posibles hacia los modelos del mundo.
\item Frans de Waal, \textit{Are We Smart Enough to Know How Smart Animals Are?}: para entender una visión de la inteligencia no antropocéntrica.
\end{itemize}

\begin{importantbox}{Juicio final}
Lo que más vale la pena llevarse de este episodio no es la dicotomía «los LLM están equivocados y los modelos del mundo tienen razón», sino un juicio más preciso: los LLM han resuelto una gran parte del mundo que la humanidad ya puso por escrito; los modelos del mundo deben resolver el mundo que la humanidad no ha escrito y que es muy difícil de escribir. Lo decisivo para la siguiente etapa de la inteligencia quizá se encuentre en el puente entre ambos.
\end{importantbox}

\end{document}
